\chapter{Study of the Existing}
\label{chap:existing}

\section{How Dependency Tracking Is Done Today}

This chapter examines the manual process used before the introduction of the
MES-X Dependency Management system. Understanding that workflow is important,
because the design choices made later in the report are direct responses to its
limitations.

In the MES X.0 environment, the starting point is usually a \textbf{user
story} in Azure DevOps. That story is often linked to several child tasks, and
each task may produce its own pull request, commits, and release references in
different repositories. Azure DevOps stores all of these artifacts, but it does
not provide a single consolidated view of their dependency chain.

In practice, a developer preparing a dependency report had to repeat the same
sequence for \emph{each} user story under analysis:

\begin{enumerate}
    \item Open the work item in Azure DevOps and note all linked child tasks and related
          items.
    \item Navigate to the \textbf{Development} panel of each item to find associated
          pull requests.
    \item Open every pull request individually, then switch to its \textbf{Commits} tab.
    \item Inspect each commit message and the list of changed files to identify which
          microservice is affected and extract any version tags present in the commit.
    \item Record the collected information---work item ID, PR number, commit hash,
          microservice name, tag---row by row in an Excel spreadsheet.
    \item Repeat the entire sequence for every child work item linked to the original
          story.
\end{enumerate}

During a typical sprint, with ten to fifteen user stories and several pull
requests per story, this work could easily occupy most of a day. Every link had
to be opened manually, and every relevant tag had to be checked one by one.

\section{Illustrative Example}

Consider a user story such as \emph{``Update shared configuration service for
plant B.''} The change may involve the configuration service itself, a dependent
API gateway, and a deployment adjustment. In practice, that single story can be
split into several child tasks, each with its own pull request and its own set
of commits.

To build a dependency view for that story, the developer must inspect each pull
request, identify the affected repository, review the associated commits, and
locate the relevant release tags. Even in a modest case, this means opening
several pull requests and reading through dozens of commits. Once this effort is
multiplied across a full sprint, the workload becomes substantial and the risk
of omission increases.

\section{Limitations and Their Consequences}

The manual process creates four major problems.

\textbf{Time consumption.} Preparing a complete dependency report for a sprint
can take from several hours to a full working day. That time is repeatedly
spent on navigation and transcription instead of implementation or review.

\textbf{High error risk.} The result depends on careful human inspection. A
single missed commit, an overlooked repository, or an incorrect tag reference
can produce an incomplete report and lead to poor release decisions.

\textbf{Poor scalability.} The MES X.0 platform already contains many
microservices, and its scope continues to grow. As more repositories are added,
the manual approach becomes less realistic and more costly.

\textbf{Lack of a shared source of truth.} Because the final result is often
stored in spreadsheets, two engineers can produce different analyses for the
same feature with no automatic way to reconcile the difference.

These limitations affect both productivity and reliability. They justify the
need for a structured, automated solution, which is formalized in the next
chapter through the system requirements.